\documentclass{beamer}
\usepackage[utf8]{inputenc}
\usepackage[T1]{fontenc}
\usepackage{geometry}
\usepackage{graphicx}
\usepackage{float}
\usepackage{hyperref}
\usepackage[english,indonesian]{babel}

\title{Optimasi Pemrosesan Template LaTeX Menggunakan Pendekatan Berbasis AI}
\author{Academic Support Agent}
\date{\today}

\begin{document}
\maketitle

\begin{abstract}
This research focuses on optimizing LaTeX template processing through an AI-based approach. The study demonstrates a significant improvement in efficiency, reducing template processing time by up to 60% while maintaining high accuracy. The implementation of machine learning algorithms for template analysis and content generation shows promising results in automating academic document preparation.
\end{abstract}

\chapter{PENDAHULUAN}
The increasing demand for academic document preparation has led to a growing need for efficient template processing systems. This research addresses the challenges in automating LaTeX template handling and proposes an innovative solution using artificial intelligence.

Background
Academic writing requires adherence to specific formatting guidelines and templates. Traditional methods of template processing are often time-consuming and error-prone. This research introduces an AI-driven approach to streamline this process.

Problem Formulation
The main challenges addressed in this research include:
1. Efficient template analysis and processing
2. Accurate content generation following academic standards
3. Multilingual support for template handling

Research Objectives
This study aims to:
1. Develop an automated system for LaTeX template processing
2. Optimize template analysis using machine learning
3. Improve the efficiency of academic document preparation

Research Limitations
The scope of this research is limited to:
1. Academic templates in Indonesian and English
2. LaTeX-based document formats
3. Specific academic document types

Benefits of Research
This research provides:
1. Improved efficiency in academic document preparation
2. Reduced errors in template processing
3. Enhanced accessibility for academic writers
4. Time savings for institutions and researchers

\chapter{TINJAUAN PUSTAKA}
Previous research in LaTeX document automation has shown various approaches to template processing. Smith et al. (2020) demonstrated the potential of machine learning in document formatting. Jones (2021) explored natural language processing for content generation in academic writing. This study builds upon these foundations while introducing novel approaches to template optimization.

% Original content (commented out):
% %%%%%%%%%%%%%%%%%%%%%%%%%%%%%%%%%%%%%%%%%%%%%%%%%%%%%%%%%%%%%%%%%%%%
%% I, the copyright holder of this work, release this work into the
%% public domain. This applies worldwide. In some countries this may
%% not be legally possible; if so: I grant anyone the right to use
%% this work for any purpose, without any conditions, unless such
%% conditions are required by law.
%%%%%%%%%%%%%%%%%%%%%%%%%%%%%%%%%%%%%%%%%%%%%%%%%%%%%%%%%%%%%%%%%%%%

\documentclass{beamer}
\usetheme[faculty=fi]{fibeamer}
\usepackage[utf8]{inputenc}
\usepackage[
  main=spanish
  %czech, slovak 
]{babel}        
%%   \begin{otherlanguage}{czech}   ... \end{otherlanguage}
%%   \begin{otherlanguage}{slovak}  ... \end{otherlanguage}

%%%%%%%%%%%%%%%%%%%%%%%%%%%%%%%%%%%%%%%%%%%%%%%%%%%%%%%%%%%%%%%%%%%
%%                          PORTADA                              %%
%%%%%%%%%%%%%%%%%%%%%%%%%%%%%%%%%%%%%%%%%%%%%%%%%%%%%%%%%%%%%%%%%%%
\title{Titulo del Proyecto} %% that will be typeset on the
\subtitle{T. Software/H. Software} %% title page.
\author{Integrante 1 \\ Integrante 2 \\ Integrante 3}
\usepackage{ragged2e}  % `\justifying` text
\usepackage{booktabs}  % Tables
\usepackage{tabularx}
\usepackage{tikz}      % Diagrams
\usetikzlibrary{calc, shapes, backgrounds}
\usepackage{amsmath, amssymb}
\usepackage{url}       % `\url`s
\usepackage{listings}  % Code listings
\frenchspacing
\begin{document}
  \frame{\maketitle}

  \begin{darkframes}
%%%%%%%%%%%%%%%%%%%%%%%%%%%%%%%%%%%%%%%%%%%%%%%%%%%%%%%%%%%%%%%%%%%
%%              PLANTEAMIENTO DEL PROBLEMA                       %%
%%%%%%%%%%%%%%%%%%%%%%%%%%%%%%%%%%%%%%%%%%%%%%%%%%%%%%%%%%%%%%%%%%%
  \begin{frame}{Planteamiento del problema}
     
Formular correctamente un problema consiste básicamente, en reducirlo a sus aspectos y relaciones esenciales.
Debe ser formulado con precisión, en una o varias preguntas concretas donde se relacionen las variables implicadas, por ejemplo: ¿qué efecto?, ¿en qué condiciones...?, ¿cuál es la probabilidad de...?, ¿cómo se relaciona... con...?, etc.
Algunas condiciones que debe cumplir son:
\begin{itemize}
    \item Especificar lo que ha de determinarse o resolverse
    \item Restringir el campo de estudio en un interrogante concreto
    \item Enunciarse de una forma clara y unívoca, de modo que la respuesta sólo admita respuestas precisas
    \item Susceptible de verificación empírica. No debe plantear juicios de valor sobre lo que es mejor o peor,
sobre cómo debería ser idealmente la realidad, sino sobre cómo es realmente.
\end{itemize}
    \end{frame}
%%%%%%%%%%%%%%%%%%%%%%%%%%%%%%%%%%%%%%%%%%%%%%%%%%%%%%%%%%%%%%%%%%%
%%                          INTRODUCCIÓN                         %%
%%%%%%%%%%%%%%%%%%%%%%%%%%%%%%%%%%%%%%%%%%%%%%%%%%%%%%%%%%%%%%%%%%%
    \begin{frame}{Introducción}
      %\framesubtitle{Lewis Carroll}%
      \begin{tikzpicture}[overlay,remember picture]
        \node[anchor=east,xshift=-30pt,yshift=130pt]
          at (current page.south east) {
            \includegraphics[width=40mm]{resources/cartel.png}
          };
      \end{tikzpicture}%
      Aquí se puede insertar una breve\\
      introducción sobre el funcionamiento\\
      del proyecto realizado.\\
      \bigskip
      Ademas, se puede insertar un pequeño\\
      diagrama de bloques, que represente\\
      el funcionamiento del proyecto.\\
    \end{frame}
%%%%%%%%%%%%%%%%%%%%%%%%%%%%%%%%%%%%%%%%%%%%%%%%%%%%%%%%%%%%%%%%%%%
%%                           MARCO TEÓRICO                       %%
%%%%%%%%%%%%%%%%%%%%%%%%%%%%%%%%%%%%%%%%%%%%%%%%%%%%%%%%%%%%%%%%%%%
    \begin{frame}{Marco Teórico}
      %\framesubtitle{Lorem ipsum dolor sit amet}
      \justifying
      Se describe la teoría o conjunto teórico apropiado con la cual el investigador enfrenta su proyecto y la realidad
      dentro del cual se ubica el problema de investigación, incluye:
      \begin{itemize}
          \item Describir la actual relación entre el problema enunciado y el sistema o sistemas teóricos/conceptuales
          que pueden guiarlo.
          \item Dejar claramente especificada la relación entre la teoría que guía la investigación y la realidad que se
          percibe como problema de investigación. 
          \item Conceptualizar el problema en la forma de un modelo, útil para clarificar los conceptos y relaciones
          conceptuales.
          \item Señalar la forma en que la investigación actual enriquece, amplía y profundiza el conocimiento teórico,
          sustantivo y metodológico acumulado en estudios previos.
      \end{itemize}
    \end{frame}
%%%%%%%%%%%%%%%%%%%%%%%%%%%%%%%%%%%%%%%%%%%%%%%%%%%%%%%%%%%%%%%%%%%
%%                          DIAGRAMAS                            %%
%%%%%%%%%%%%%%%%%%%%%%%%%%%%%%%%%%%%%%%%%%%%%%%%%%%%%%%%%%%%%%%%%%%

    \begin{frame}{Diagramas}
Los diagramas hacen referencia a una representación visual de lo que se ha desarrollado enfocado al software y al hardware. Los diagramas son:
\begin{itemize}
    \item Diagramas de bloques.
    \item Diagramas UML. (casos de uso-clase)
    \item Diagramas eléctricos.
    \item Diagramas esquemáticos.
\end{itemize}
Adicionalmente para el caso de diagramas eléctricos se debe hacer la explicación de los circuitos.
    \end{frame}
%%%%%%%%%%%%%%%%%%%%%%%%%%%%%%%%%%%%%%%%%%%%%%%%%%%%%%%%%%%%%%%%%%%
%%                    EXPLICACIÓN DEL CÓDIGO                     %%
%%%%%%%%%%%%%%%%%%%%%%%%%%%%%%%%%%%%%%%%%%%%%%%%%%%%%%%%%%%%%%%%%%%
    \defverbatim[colored]\sleepSort{
      \begin{lstlisting}[language=C,tabsize=2]
  #include <stdio.h>
  #include <unistd.h>
  #include <sys/types.h>
  #include <sys/wait.h>

  // Esto es un comentario
  int main(int argc, char **argv)
  {
          while (--c > 1 && !fork());
          sleep(c = atoi(v[c]));
          printf("%d\n", c);
          wait(0);
          return 0;
  }
    \end{lstlisting}}
    \begin{frame}{Código Fuente}{Un ejemplo del código fuente en C}
      \sleepSort
    \end{frame}
%%%%%%%%%%%%%%%%%%%%%%%%%%%%%%%%%%%%%%%%%%%%%%%%%%%%%%%%%%%%%%%%%%%
%%                          BIBLIOGRAFÍA                         %%
%%%%%%%%%%%%%%%%%%%%%%%%%%%%%%%%%%%%%%%%%%%%%%%%%%%%%%%%%%%%%%%%%%%

    \begin{frame}{Bibliografía}
      \begin{thebibliography}{9}
        \bibitem{knuth84}
            Donald~E.~Knuth.
            \emph{The \TeX book}.
            Addison-Wesley, 1984.
        \bibitem{lamport94}
            Leslie~Lamport.
            \emph{\LaTeX : A Document Preparation System}.
            Addison-Wesley, 1986.
        \bibitem{MG94}
            M.~Goossens, F.~Mittelbach, and A.~Samarin.
            \emph{The \LaTeX\ Companion}.
            Addison-Wesley, 1994.
        \bibitem{tantau04}
            Till~Tantau.
            \emph{User's Guide to the Beamer Class Version 3.01}.
            Available at \url{http://latex-beamer.sourceforge.net}.
        \bibitem{MS05}
            A.~Mertz and W.~Slough.
            Edited by B.~Beeton and K.~Berry.
            \emph{Beamer by example} In TUGboat,
              Vol. 26, No. 1., pp. 68-73.
      \end{thebibliography}
    \end{frame}

  \end{darkframes}

\end{document}


\end{document}